% ==============================================================================
% SLIDE 8: HEAVY INFRASTRUCTURE GEOTECHNICS: CABLE CAR
% ==============================================================================
\begin{frame}{Heavy Infrastructure: Anchoring the Santorini Cable Car}
  \begin{columns}[T, onlytextwidth]
    % Left Column: Geotechnical Diagram
    \begin{column}{0.48\textwidth}
      \slidephoto{cable-car.jpg}{Geotechnical anchoring: Consolidation grouting, 30\,m pre-stressed rock anchors, and micropile groups}
    \end{column}
    
    % Right Column: Bullet Points
    \begin{column}{0.48\textwidth}
      \textbf{\color{aegeanblue}Deep Rock Stabilization (Doppelmayr 1982):}
      \vspace{0.2em}
      \begin{itemize}
        \setlength{\itemsep}{0.3em}
        \item \textbf{The Engineering Challenge:} Anchoring dynamic pylon loads (220\,m vertical rise, 1,200 passengers/hour) into fractured, crumbly volcanic rock.
        \item \textbf{High-Pressure Grout Injection:} Cementitious grout injected into fractured zones glues loose blocky lava together, drastically increasing the Young's modulus of the rock mass.
        \item \textbf{Pre-Stressed Rock Anchors:} Steel tendons drilled 15--30\,m deep penetrate far past active cliff slip planes into competent bedrock, hydraulically clamping surface rock under compression.
        \item \textbf{Distributed Micropiling:} Vertical micropiles dissipate concentrated footing weights through skin friction along deep columns.
      \end{itemize}
    \end{column}
  \end{columns}
\end{frame}
